\section{Conclusion}
This literature review has examined the limitations of conventional front-side PDNs, the advantages of back-side PDNs, and the challenges introduced by these architectures in advanced technology nodes. Scaling beyond the 5nm range, increased interconnect resistance,  reduced voltage margins and routing congestion have made traditional front-side PDNs increasingly inefficient and difficult to optimise. 

The reviewed studies demonstrate that back-side PDNs provide substantial improvements in power integrity and design efficiency. By relocating the power network to the backside of the wafer, this architecture decouples power and signal lines, thereby increasing routing efficiency and better utilizing front-end metal layers. Major reductions in IR drop have been demonstrated, with reported improvements of up to 85\% compared to front-side PDNs designs. Furthermore, back-side implementations have shown clear gains in power, performance, and area, including frequency improvements of 6\% and core area reductions of about 20\% and upwards depending on the design and technology scenario.

Despite these advantages, back-side PDNs introduce new challenges. The relocation of the power network affects heat dissipation, resulting in higher thermal resistance and hotspots with a 45\% increase in temperature. Furthermore, integration of additional architectural features such as Buried Power Rails and nano-TSVs introduces additional parasitic effects that need to be carefully managed.

In conclusion, back-side PDNs demonstrate a significant evolution in power delivery architecture that addresses limitations of conventional front-side PDNs. Although back-side PDNs introduce new challenges in thermal management and fabrication, the architecture's improvements in IR drop, performance and cell area make it a strong candidate for technology nodes in the future. While further research in these limitations is necessitated, current evidence indicated that back-side PDNs will have a significant role in sub-5nm CMOS design.